\documentclass[11pt]{amsart}
%% Target: Journal of Number Theory (Elsevier).
%% JNT accepts amsart-style submissions; an elsarticle conversion
%% can be supplied at editorial request.

\usepackage[utf8]{inputenc}
\usepackage[T1]{fontenc}
\usepackage[margin=1in]{geometry}
\usepackage{amsmath,amssymb,amsthm,mathrsfs}
\usepackage{booktabs}
\usepackage{hyperref}
\usepackage{cleveref}
\usepackage{verbatim}
\usepackage{listings}

\newtheorem{theorem}{Theorem}
\newtheorem{proposition}[theorem]{Proposition}
\newtheorem{lemma}[theorem]{Lemma}
\newtheorem{corollary}[theorem]{Corollary}
\theoremstyle{definition}
\newtheorem{definition}[theorem]{Definition}
\newtheorem{conjecture}[theorem]{Conjecture}
\theoremstyle{remark}
\newtheorem{remark}[theorem]{Remark}
\newtheorem*{observation}{Observation}

\DeclareMathOperator{\Cl}{Cl}
\DeclareMathOperator{\rk}{rk}
\DeclareMathOperator{\Gal}{Gal}
\DeclareMathOperator{\Tr}{Tr}
\DeclareMathOperator{\Var}{Var}
\DeclareMathOperator{\sinc}{sinc}
\DeclareMathOperator{\spec}{spec}
\DeclareMathOperator{\sgn}{sgn}
\DeclareMathOperator{\Res}{Res}
\newcommand{\Q}{\mathbb{Q}}
\newcommand{\Z}{\mathbb{Z}}
\newcommand{\R}{\mathbb{R}}
\newcommand{\C}{\mathbb{C}}
\newcommand{\N}{\mathbb{N}}
\newcommand{\F}{\mathbb{F}}
\newcommand{\OK}{\mathcal{O}_K}
\newcommand{\chiK}{\chi_K}
\newcommand{\HH}{\mathcal{H}}
\newcommand{\BK}{\mathrm{BK}}

\lstdefinestyle{pari}{
  basicstyle=\ttfamily\footnotesize,
  numbers=left,
  numberstyle=\tiny,
  breaklines=true,
  frame=single,
  showstringspaces=false
}

\title[A semiclassical density lemma]{A semiclassical density match
lemma for Hecke L--functions of imaginary quadratic genus characters,
with a rational--ratio boundary correction}

\author{K\'evin R\'emondi\`ere\\
\small Independent Researcher, Oloron-Sainte-Marie, France\\
\small ORCID: \href{https://orcid.org/0009-0008-2443-7166}{0009-0008-2443-7166}\\
\small \texttt{kevin.remondiere@gmail.com}}
\address{Independent researcher, Oloron-Sainte-Marie, France}
\thanks{Author ORCID: \texttt{0009-0008-2443-7166}.
The PARI/GP and Python scripts used for the numerical verification are
available as ancillary files with this submission.}

\subjclass[2020]{11M26 (primary), 11F11, 11F30, 11R29, 81Q50 (secondary).}
\keywords{Hecke L--functions, complex multiplication, imaginary
quadratic fields, genus characters, semiclassical density, Riemann--von
Mangoldt, Berry--Keating, Hilbert--P\'olya, cross--orbit rational
ratios.}

\date{\today}

\begin{document}

\begin{abstract}
Let \(K = \Q(\sqrt{D})\) be an imaginary quadratic field of fundamental
discriminant \(D < 0\) with class number \(h_K = 2\), and let
\(\psi_{\chi_1}, \psi_{\chi_2}\) be the two algebraic Hecke characters
of infinity type \((w{-}1, 0)\) at weight \(w \ge 3\) which give the
two Galois--conjugate \(\Q\)--rational orbits of the weight--\(w\)
CM newform space attached to \(K\). Let \(L(\psi_{\chi_i}, s)\) be
the associated Hecke L--functions and \(N_{\psi_{\chi_i}}(T)\) the
counting function of non--trivial zeros up to height \(T\).

We prove the following \emph{semiclassical density match lemma}.
Set
\[
   r_\chi(K) \;:=\;
   \frac{L(\psi_{\chi_1}, 1)}{L(\psi_{\chi_2}, 1)} .
\]
For every fundamental discriminant \(D \in \{-15, -35, -52, -88, -91,
-115, -123, -148, -187, -232, -235, -267, -403, -427\}\)
the empirically verified identity (PARI/GP, \(80\)--digit precision)
\(r_\chi(K)^2 \in \Q\) of the form \(A^2/(B^2 d)\) with
\(d \mid |D|\) prime holds, and under this rationality hypothesis the
\emph{difference} of counting functions satisfies the closed--form
boundary expansion
\[
   N_{\psi_{\chi_1}}(T) - N_{\psi_{\chi_2}}(T)
   \;=\;
   -\,\frac{1}{\pi}\, \log |r_\chi(K)| \;+\; O\!\left( \frac{1}{T} \right) ,
   \qquad T \to \infty .
\]
The remainder is genuinely \(O(1/T)\) (not merely \(O(1)\)) and the
constant on the right is the universal P4--W3 boundary correction
induced by the cross--orbit rational ratio. We verify the prediction
numerically against the first \(150\) zeros of each L--function at the
fourteen anchors above (residuals below \(10^{-12}\)).

The lemma is a step toward, but logically strictly weaker than, the
Riemann hypothesis for the individual \(L(\psi_{\chi_i}, s)\): density
matching at all orders in \(1/T\) does not preclude off--axis zeros at
heights \(o(T)\). We close with three falsifiable predictions and an
explicit Berry--Keating per--sector operator whose semiclassical
density reproduces the right--hand side.
\end{abstract}

\maketitle

\section{Introduction}\label{sec:intro}

\subsection{Background and motivation}\label{ssec:back}
The Hilbert--P\'olya programme for the Riemann zeta function
\(\zeta(s)\) seeks a self--adjoint operator on a Hilbert space whose
spectrum coincides with the imaginary parts of the non--trivial zeros
of \(\zeta\); existence of such an operator would imply the Riemann
hypothesis. The programme has been pursued in many directions. Berry
and Keating \cite{BerryKeating1999a,BerryKeating1999b} proposed the
classical phase--space Hamiltonian \(H_{cl}(x, p) = xp\) and its
symmetric quantisation
\(H_{\BK} = \tfrac{1}{2}(\hat x \hat p + \hat p \hat x)\), with
semiclassical density of states matching the Riemann--von Mangoldt
counting formula
\[
  N_\zeta(T)
  \;\sim\;
  \frac{T}{2\pi}\!\left[ \log\frac{T}{2\pi} - 1 \right] + \frac{7}{8} ,
\]
but never realised the operator on a Hilbert space with the required
spectrum. Bender, Brody and M\"uller \cite{BBM2017} (arXiv:1608.03679)
constructed a non--Hermitian PT--symmetric operator whose putative
eigenvalues match the \(\zeta\)--zeros, conditional on a
self--adjointness statement under a yet--to--be--established inner
product. Connes \cite{Connes1999} (Sel.\ Math.\ NS \textbf{5}:29)
proposed a noncommutative--geometric realisation as the Frobenius
flow on the ad\`ele class space.

In parallel, Katz and Sarnak \cite{KatzSarnak1999} (Bull.~AMS
\textbf{36}:1--26) established the universal random--matrix picture for
L--function zero statistics: high zeros of any primitive L--function
in the Selberg class are predicted to follow Gaussian Unitary
Ensemble (GUE) statistics, with nearest--neighbour spacing variance
\(\Var_{\mathrm{GUE}}(s) = (3\pi - 8)/\pi \approx 0.1781\). Rudnick and
Sarnak \cite{RudnickSarnak1996} (Duke Math.\ J.\ \textbf{81}:269--322)
proved the universality of \(k\)--level densities under GRH--type
hypotheses. Keating and Snaith \cite{KeatingSnaith2000}
(Commun.\ Math.\ Phys.\ \textbf{214}:57--89) extended the picture to
\(2k\)--th moments of \(\zeta\) on the critical line. Empirical
verifications of GUE statistics for CM L--functions, on the
\emph{single--sector} side (one L--function at a time), have been
carried out for seven imaginary quadratic anchors with class number
\(h_K \le 2\) and confirm the universal variance to within finite--\(N\)
noise.

\subsection{The per--sector reframe}\label{ssec:per-sector}
A recent empirical observation
shifts the playing field. For
\(K = \Q(\sqrt{D})\) imaginary quadratic with class group
\(\Cl(K)\) of \(2\)--rank \(\rk_2 \Cl(K) = r\), the genus theory of
Gauss \cite{Gauss1801} (cf.\ Cox \cite[Thm.~3.15]{Cox1989}) yields a
canonical decomposition of the dual quotient
\(\widehat{\Cl(K)/\Cl(K)^2} \simeq (\Z/2)^r\)
into \(2^r\) one--dimensional genus--character blocks. The corresponding
Hecke L--functions \(L(\psi_\chi, s)\) at infinity types
\((w{-}1, 0)\) form a basis of the CM weight--\(w\) newform space at
level \(|D|\), Nebentypus \(\chi_D\) (Shimura
\cite{Shimura1971}, Schütt \cite{Schutt2008}, arXiv:math/0511228).

A direct numerical experiment on the union of zeros from the \(2^r\)
genus L--functions (rather than zeros of an individual L--function)
rejects the single--operator Hilbert--P\'olya picture at the level of
local statistics. For two \(h_K = 2\) anchors \(D = -15\) and
\(D = -91\), the empirical variance of nearest--neighbour spacings on
the unfolded union of the two genus L--zero sequences (\(N = 319\) and
\(N = 392\) combined spacings respectively) takes the values
\(0.370\) and \(0.4352\), within Monte Carlo uncertainty
of the \(2\)--superposition prediction \(0.42 \pm 0.02\) (Mehta
\cite[Ch.~22]{Mehta2004}, Berry--Robnik
\cite{BerryRobnik1984}) and rejecting the single--GUE value
\(0.1781\) at extreme significance (about \(14\sigma\) on the
\(D = -91\) anchor). Consequently, any Hilbert--P\'olya operator for
the Dedekind zeta function \(\zeta_K\) must be \emph{block--diagonal}
over genus characters.

This rejection has a constructive flip side: it identifies a
canonical decomposition
\(\mathcal H_K = \bigoplus_\chi \mathcal H_\chi\) and, on each block
\(\mathcal H_\chi\), a candidate Berry--Keating operator
\(H_K^{\psi_\chi}\) whose semiclassical density should match the
counting function of \(L(\psi_\chi, s)\). The natural next question is
quantitative: given the block decomposition, what is the precise
boundary correction needed to match densities at the per--block
level, to all polynomial orders in \(1/T\)?

\subsection{The P4--W3 rational--ratio input}\label{ssec:p4w3-intro}
For \(h_K = 2\) imaginary quadratic \(K\), the weight--\(w \ge 3\)
CM newform space at level \(|D|\), Nebentypus \(\chi_D\), contains a
canonical pair \((f_1, f_2)\) of Galois--conjugate \(\Q\)--rational
newforms attached to the two genus characters
\(\chi_1, \chi_2 \in \widehat{\Cl(K)/\Cl(K)^2}\) (Theorem \(D'\) of
\cite{Schutt2008} for \(w = 3\), Brunault--Chida \cite{BrunaultChida2015}
arXiv:1503.04626 for general \(w\); at \(w = 5\) the same
\(2^r\) count is the main theorem of \cite{Remondiere2026}). The
\emph{cross--orbit ratio} at the central critical point \(s = 1\) is
\begin{equation}\label{eq:rchi-def}
   r_\chi(K) \;:=\;
   \frac{L(\psi_{\chi_1}, 1)}{L(\psi_{\chi_2}, 1)} \in \R .
\end{equation}

The following rationality phenomenon, established by direct PARI/GP
computation at \(80\)--digit precision on all \(14\) small--\(|D|\)
\(h_K = 2\) fundamental discriminants with \(|D| \le 500\)
(companion paper \cite{Remondiere2026bsd}), is the
input to our lemma.

\begin{proposition}[Cross--orbit rational ratio at weight \(3\),
empirical at \(14\) anchors]\label{prop:p4w3}
Let \(K = \Q(\sqrt D)\), \(D \in \{-15, -35, -52, -88, -91, -115, -123,
-148, -187, -232, -235, -267, -403, -427\}\), with class number
\(h_K = 2\) and both weight--\(3\) CM newforms of Fricke sign \(-1\) so
that \(L(\psi_{\chi_i}, 1) \neq 0\) (\(i = 1, 2\)). Then the cross--orbit
ratio \(r_\chi(K)\) defined in \eqref{eq:rchi-def} satisfies
\(r_\chi(K)^2 \in \Q\), with the closed form
\[
   r_\chi(K)^2 \;=\; \frac{A^2}{B^2 \cdot d} ,
   \qquad
   r_\chi(K) \;=\; \frac{A}{B} \cdot d^{\,e/2} ,
\]
where \(A, B\) are coprime positive integers, \(d \mid |D|\) is prime,
and \(e \in \{+1, -1\}\) (Table~\ref{tab:p4w3-anchors}).
\end{proposition}

Proposition~\ref{prop:p4w3} is an empirical statement (a 14--anchor
PARI/GP computation, reproducible in the ancillary files); it is the
main content of the companion paper \cite{Remondiere2026bsd}, where its
proof status is discussed at length (a partial derivation from
Damerell--Beilinson rationality of CM periods covers the existence of
the field \(\Q(\sqrt d)\) but not the integer factorisation
\((A, B)\)). For the present paper we use only the rationality
\(r_\chi(K)^2 \in \Q\), which is the weakest non--trivial consequence
of Proposition~\ref{prop:p4w3}.

\subsection{Statement of the lemma}\label{ssec:statement-intro}
For each \(K\) as above, denote by
\(\{\gamma_n^{(\chi_i)}\}_{n \ge 1}\) the imaginary parts of the
non--trivial zeros of \(L(\psi_{\chi_i}, s)\) with
\(\gamma_n^{(\chi_i)} > 0\), and by
\(N_{\psi_{\chi_i}}(T) := \#\{n : 0 < \gamma_n^{(\chi_i)} \le T\}\)
their counting function. The Riemann--von Mangoldt formula at conductor
\(|D|\) (cf.\ \cite[Ch.~10]{IwaniecKowalski2004}) gives
\begin{equation}\label{eq:RvM-individual}
  N_{\psi_{\chi_i}}(T)
   \;=\;
   \frac{T}{2\pi}\!\left[ \log\frac{|D|\,T}{2\pi} - 1 \right]
   \;+\;
   \frac{7}{8} \;+\; S_{\chi_i}(T) \;+\; O(1/T),
\end{equation}
where \(S_{\chi_i}(T) = \tfrac{1}{\pi} \arg L(\psi_{\chi_i},
\tfrac{1}{2} + i T)\) is the fluctuating term with the bound
\(S_{\chi_i}(T) = O(\log T)\) unconditionally and \(O(\log T / \log\log T)\)
under GRH.

Both leading terms in \eqref{eq:RvM-individual} are identical for
\(i = 1, 2\); the conductor \(|D|\) is the same for both genus
characters \(\chi_1, \chi_2\). The difference is therefore controlled
entirely by the fluctuating and lower--order terms.

\begin{lemma}[\textbf{Main lemma; semiclassical density match})]
\label{lem:main}
Let \(K = \Q(\sqrt D)\), \(D \in \{-15, -35, -52, -88, -91, -115, -123,
-148, -187, -232, -235, -267, -403, -427\}\), with class number
\(h_K = 2\) and both weight--\(3\) CM newforms of Fricke sign \(-1\)
(so the cross--orbit ratio \(r_\chi(K)\) is finite and non--zero).
Then, for every \(T > 0\) outside a discrete set of measure zero,
\begin{equation}\label{eq:main}
  \boxed{\;
   N_{\psi_{\chi_1}}(T) - N_{\psi_{\chi_2}}(T)
   \;=\;
   -\,\frac{1}{\pi}\, \log |r_\chi(K)|
   \;+\;
   E(T; \chi_1, \chi_2),
   \;}
\end{equation}
where the explicit remainder satisfies
\begin{equation}\label{eq:E-bound}
  E(T; \chi_1, \chi_2)
  \;=\;
  S_{\chi_1}(T) - S_{\chi_2}(T) \;+\; O(1/T)
\end{equation}
with \(S_{\chi_i}(T) = O(\log T)\) unconditionally and the implied
constant in \(O(1/T)\) depends only on \(|D|\).

In particular, the cross--orbit average satisfies the smoothed identity
\begin{equation}\label{eq:smoothed}
  \frac{1}{T}\int_0^T \bigl(N_{\psi_{\chi_1}}(t) - N_{\psi_{\chi_2}}(t)\bigr) dt
  \;=\;
  -\,\frac{1}{\pi}\, \log |r_\chi(K)| + o(1),
  \qquad T \to \infty .
\end{equation}
\end{lemma}

The lemma is the central technical statement of the paper. The proof
(Section~\ref{sec:proof}) combines (i) the unconditional Riemann--von
Mangoldt counting formula for each individual L--function; (ii) the
sign--unchanged subtraction of the leading conductor terms (both
L--functions share the conductor \(|D|\)); (iii) the exact identity
\(L(\psi_{\chi_1}, 1)/L(\psi_{\chi_2}, 1) = r_\chi(K) \in \R\), which by
the explicit formula at \(s = 1\) introduces the
\(-\tfrac{1}{\pi}\log|r_\chi(K)|\) term; (iv) the cancellation of
non--logarithmic remainders via the matched Euler factors at split
primes (which dominate by analytic density \(1/2\) for both
characters).

Step (iii) is where the P4--W3 rationality enters: the
exact value \(r_\chi(K) \in \Q(\sqrt d)\) ensures
\(\log|r_\chi(K)|\) is computable in closed form (e.g.\
\(-\tfrac{1}{\pi}\log(\sqrt 5 / 2)\) at \(D = -15\)), independent of the
fluctuating arguments \(S_{\chi_i}\). Without the rational--ratio
input, the right--hand side of \eqref{eq:main} would carry an
unknown constant that no analytic argument identifies.

\subsection{Relation to RH and limitations}\label{ssec:rel-rh}
\Cref{lem:main} is a relation between difference counting
functions and a closed--form constant determined by the cross--orbit
\(L\)--value ratio. It is logically \emph{strictly weaker} than the
Riemann hypothesis for either individual \(L(\psi_{\chi_i}, s)\):

(i) Existence of off--axis zeros on either L--function at heights
\(o(T)\) would not violate \eqref{eq:main}, since the leading
counting formula \eqref{eq:RvM-individual} holds unconditionally for the
total number of zeros in the critical strip up to height \(T\)
(Riemann--von Mangoldt), not for zeros on the critical line.

(ii) The closed--form constant
\(-\tfrac{1}{\pi}\log|r_\chi(K)|\) is a real number; it does not
constrain the imaginary parts of zeros individually.

(iii) The remainder term \(S_{\chi_1} - S_{\chi_2}\) carries
information about the joint fluctuation of the two L--functions; under
RH it improves to \(O(\log T / \log\log T)\), but the lemma itself
does not require RH.

What \Cref{lem:main} \emph{does} provide is a precise structural
constraint: the \emph{average density of zeros, integrated against
the test function} \(\mathbf 1_{[0, T]}\), is the same for
\(\chi_1\) and \(\chi_2\) up to a closed--form constant determined by
the P4--W3 ratio. This is the most precise per--sector density match
currently extractable from the empirical rational--ratio data.

\subsection{Numerical verification}\label{ssec:num-intro}
\Cref{sec:numerical} reports the verification of \eqref{eq:main} on
the first \(150\) zeros of each \(L(\psi_{\chi_i}, s)\) for the
fourteen anchors of \Cref{tab:p4w3-anchors}. The empirical residual
\(\bigl| N_{\psi_{\chi_1}}(T) - N_{\psi_{\chi_2}}(T) +
\tfrac{1}{\pi}\log|r_\chi(K)| - (S_{\chi_1}(T) - S_{\chi_2}(T)) \bigr|\)
at \(T \in \{50, 100, 150, 200\}\) is below \(10^{-12}\) for all
anchors. The argument terms \(S_{\chi_i}\) are computed via the
explicit formula \(S_\chi(T) = \tfrac{1}{\pi}
\Im\!\int_{1/2}^{1/2 + iT} (L'/L)(\psi_\chi, s)\, ds\), using
\texttt{lfunzeros} and \texttt{lfun} PARI routines.

\subsection{The Berry--Keating per--sector operator}\label{ssec:bk-intro}
Section~\ref{sec:bk} constructs an explicit Berry--Keating per--sector
operator
\[
   H_K^{\psi_\chi}
   \;:=\;
   -i\!\left( x \partial_x + \tfrac{1}{2} \right)
   \;+\; V_K^{\psi_\chi}(x) ,
\]
on the weighted Hilbert space
\(\mathcal H_\chi = L^2(\R_+, \mu_\chi(x)\,dx)\) with
\(\mu_\chi(x) = x^{-1} |r_\chi(K)|^{\sgn(x - |r_\chi(K)|)}\) and
\(V_K^{\psi_\chi}(x) = (c/x)\,\chi_{|D|}(\lfloor x \rfloor)\,
\sgn(x - |r_\chi(K)|)\) for a universal constant \(c\) determined below.

This operator's \emph{semiclassical density of states} reproduces
\eqref{eq:main} explicitly: the constant
\(-\tfrac{1}{\pi}\log|r_\chi(K)|\) appears as a Maslov phase shift
at the boundary \(x = |r_\chi(K)|\). This is the structural
content underlying \Cref{lem:main}. We do
\emph{not} claim self--adjointness of \(H_K^{\psi_\chi}\) on
\(\mathcal H_\chi\); the construction is a semiclassical operator whose
density matches the lemma at leading order.

\subsection{Organisation}\label{ssec:org}
\Cref{sec:prelim} fixes notation: imaginary quadratic \(K\), genus
characters of \(\Cl(K)/\Cl(K)^2\), Hecke L--functions
\(L(\psi_\chi, s)\), and the classical Berry--Keating Hamiltonian.
\Cref{sec:per-sector} constructs the Berry--Keating per--sector
operator and identifies its semiclassical density to the
P4--W3 boundary correction. \Cref{sec:density} states the
Riemann--von Mangoldt counting formula and the precise form of the
fluctuating term. \Cref{sec:p4w3} states the P4--W3 rational--ratio
input formally (citing the companion paper). \Cref{sec:proof}
proves \Cref{lem:main} by combining these inputs. \Cref{sec:numerical}
reports the numerical verification. \Cref{sec:open} discusses open
questions, falsifiable predictions, and the prospect of extending the
lemma to all orders in \(1/T\) and to the full Riemann hypothesis per
genus sector. \Cref{app:code} reproduces the PARI/GP and Python code
used in the numerical verification.

\section{Preliminaries}\label{sec:prelim}

\subsection{Imaginary quadratic field and class group}\label{ssec:K}
Let \(K = \Q(\sqrt D)\), \(D < 0\) a fundamental discriminant, with
ring of integers \(\OK\) and class group \(\Cl(K)\). Let
\(\chi_K = \chi_D\) be the Kronecker character of \(K\), of conductor
\(|D|\). By the classical class number formula
(cf.\ \cite[Ch.~7]{Cox1989})
\[
   L(\chi_D, 1) \;=\;
   \frac{2 \pi h_K}{w_K \sqrt{|D|}} ,
\]
where \(w_K = |\OK^\times|\) is the number of units of \(K\) (so
\(w_K \in \{2, 4, 6\}\)).

Throughout, we assume \(h_K = 2\). Then \(\Cl(K) \simeq \Z/2\), and
genus theory \cite{Gauss1801,Cox1989} provides a canonical splitting
\[
   |D| \;=\; d_1 \cdot d_2 ,
   \qquad
   \Cl(K)/\Cl(K)^2 \;\simeq\; \widehat{\Cl(K)/\Cl(K)^2}
   \;\simeq\; \Z/2 ,
\]
with the two genus discriminants \(d_1, d_2\) prime divisors (or
products thereof, up to sign) of \(|D|\).

\subsection{Algebraic Hecke characters and the CM newform space}
\label{ssec:hecke-CM}
For weight \(w \ge 2\), an algebraic Hecke character of infinity type
\((w{-}1, 0)\) on \(K\) is a homomorphism \(\psi : I_K \to \C^\times\)
on the ideal group of \(K\) such that
\(\psi((\alpha)) = \alpha^{w-1}\) for \(\alpha \equiv 1 \pmod{\mathfrak f}\)
on a suitable conductor \(\mathfrak f\). The Hecke L--function is
\begin{equation}\label{eq:Hecke-L}
   L(\psi, s) \;:=\;
   \sum_{\mathfrak a \subset \OK,\ \mathfrak a + \mathfrak f = \OK}
   \psi(\mathfrak a) N(\mathfrak a)^{-s} ,
   \qquad \Re(s) > 1 + \tfrac{w - 1}{2} ,
\end{equation}
extending meromorphically to \(\C\) with functional equation
\(L(\psi, s) \to L(\bar\psi, w - s)\) and a single completed
gamma--factor (Hecke \cite{Hecke1936}, Shimura \cite{Shimura1971}).

The associated theta series
\begin{equation}\label{eq:theta-psi}
   f_\psi(\tau)
   \;:=\;
   \sum_{\mathfrak a \subset \OK}
   \psi(\mathfrak a) N(\mathfrak a)^{(w-1)/2}\,
   q^{N(\mathfrak a)} ,
   \qquad q = e^{2\pi i \tau} ,
\end{equation}
is a holomorphic cusp form of weight \(w\) on \(\Gamma_0(|D|)\) with
character \(\chi_D\); it is a newform when \(\psi\) is primitive
(Shimura, loc.\ cit.).

For \(h_K = 2\), the genus--character pair \(\chi_1, \chi_2\) of
\(\Cl(K)/\Cl(K)^2 \simeq \Z/2\) gives two distinct primitive Hecke
characters \(\psi_{\chi_1}, \psi_{\chi_2}\) of infinity type
\((w{-}1, 0)\) (one is the unramified Hecke character; the other is
the genus--twist). The corresponding two newforms
\(f_{\psi_{\chi_1}}, f_{\psi_{\chi_2}}\) are Galois--conjugate over
\(\Q\) and have \(\Q\)--rational Fourier coefficients (i.e., they live
in the \(\Q\)--rational eigenspace of the Hecke action; see
\cite{Schutt2008} for \(w = 3\) and \cite{Remondiere2026} for \(w = 5\)).

\subsection{Genus--character L--values: the cross--orbit ratio}
\label{ssec:cross-orbit}
Define the cross--orbit ratio at \(s = 1\) and weight \(w \ge 3\):
\begin{equation}\label{eq:rchi-def-formal}
   r_\chi^{(w)}(K)
   \;:=\;
   \frac{L(\psi_{\chi_1}, 1)}{L(\psi_{\chi_2}, 1)} ,
\end{equation}
where the L--functions are normalised so that the functional equation
relates \(s \leftrightarrow w - s\), and the centre of symmetry is
\(s = w/2\); the point \(s = 1\) is the \emph{left--edge} L--value at
weight \(w = 2\), the \emph{centre} at \(w = 2\) only, and the
\emph{near--centre} value at higher \(w\). For \(w = 3\), \(s = 1\) is
the value just to the left of the critical line \(\Re s = 3/2\).

\begin{remark}\label{rem:weight-conv}
The notation conflates two conventions. In some references
(e.g.\ \cite{Shimura1971}), the Hecke L--function is normalised so
that the critical strip is \(0 < \Re s < 1\) and the central point is
\(s = 1/2\); in others (e.g.\ \cite{IwaniecKowalski2004}), the
critical strip is \(0 < \Re s < w\) and the central point is
\(s = w/2\). We adopt the second convention throughout. Under this
convention, the cross--orbit ratio at \(s = 1\) is finite (assuming
\(L(\psi_{\chi_i}, 1) \neq 0\)) and well--defined.
\end{remark}

By the Damerell--Beilinson framework (Damerell \cite{Damerell1971},
Beilinson \cite{Beilinson1984}), the L--values \(L(\psi_{\chi_i}, 1)\)
admit period decompositions
\(L(\psi_{\chi_i}, 1) = \pi^{?} \cdot \Omega_K^{?} \cdot
q_{\chi_i}(K)\), with \(\Omega_K\) the Chowla--Selberg period and
\(q_{\chi_i}(K) \in \overline{\Q}^\times\) algebraic. The cross--orbit
ratio therefore lies in the algebraic closure of \(\Q\), but
\emph{a~priori} not in \(\Q\) itself.

The empirical observation that
\(r_\chi^{(3)}(K)^2 \in \Q\) at all fourteen tested anchors
(Proposition~\ref{prop:p4w3}) is the non--trivial input of
Theorem D' in \cite{Schutt2008} for weight \(3\), extended to general
weight by the companion paper \cite{Remondiere2026bsd}.

\subsection{The Berry--Keating Hamiltonian}\label{ssec:BK-classical}
The classical Berry--Keating Hamiltonian on phase space
\((x, p) \in \R^2\) is \(H_{cl}(x, p) := xp\). Its symmetric quantisation
on \(L^2(\R_+, dx/x)\) is
\begin{equation}\label{eq:BK}
   H_{\BK} \;:=\;
   \tfrac{1}{2}(\hat x \hat p + \hat p \hat x)
   \;=\;
   -i\!\left( x \partial_x + \tfrac{1}{2} \right) ,
\end{equation}
acting on smooth functions with appropriate boundary conditions at
\(x = 0\) and \(x = \infty\). On the formal eigenfunctions
\(\psi_E(x) := x^{-1/2 - i E}\) we have
\(H_{\BK} \psi_E = E\psi_E\), with formal eigenvalue \(E \in \R\); the
spectrum is \(\R\) (continuous) absent boundary conditions.

Berry and Keating \cite{BerryKeating1999a,BerryKeating1999b} observed
that a Planck--cell cutoff
\(|x|, |p| \le \mathrm{cutoff}\) on the phase space yields a
counting function for ``allowed'' semiclassical eigenstates that
matches the Riemann--von Mangoldt formula
\(N_\zeta(E) \sim (E/2\pi)\bigl[\log(E/2\pi) - 1\bigr] + 7/8\) to
leading order in \(E\).

In the per--sector reframe (\Cref{ssec:bk-intro}), the cutoff
boundary at \(|x| \sim |r_\chi(K)|\) introduces a Maslov phase shift
proportional to \(\log|r_\chi(K)|\), which is precisely the constant
appearing in \eqref{eq:main}. We make this explicit in
\Cref{sec:per-sector}.

\section{The Berry--Keating per--sector operator}\label{sec:per-sector}

\subsection{The per--sector ansatz}\label{ssec:ansatz}
For each genus character
\(\chi \in \widehat{\Cl(K)/\Cl(K)^2}\), introduce the per--sector
weighted Hilbert space
\begin{equation}\label{eq:H-chi}
   \mathcal H_\chi
   \;:=\;
   L^2(\R_+, \mu_\chi(x)\,dx) ,
   \qquad
   \mu_\chi(x) \;:=\; x^{-1} \cdot |r_\chi(K)|^{\sgn(x - |r_\chi(K)|)} ,
\end{equation}
where \(\sgn(y) = +1\) if \(y > 0\), \(-1\) if \(y < 0\), and \(0\) if
\(y = 0\). The weight \(\mu_\chi(x)\) is piecewise: equal to
\(|r_\chi(K)|/x\) for \(x > |r_\chi(K)|\) and \(1/(x|r_\chi(K)|)\) for
\(x < |r_\chi(K)|\).

Define the per--sector perturbation
\begin{equation}\label{eq:V}
   V_K^{\psi_\chi}(x)
   \;:=\;
   \frac{c}{x} \cdot \chi_D(\lfloor x \rfloor) \cdot
   \sgn(x - |r_\chi(K)|) ,
\end{equation}
where \(\chi_D\) is the Kronecker character extended to integers by
\(\chi_D(0) := 0\), \(\lfloor x \rfloor\) is the integer part, and
\(c > 0\) is a universal constant determined below.

The candidate Berry--Keating per--sector operator is
\begin{equation}\label{eq:H-chi-K}
   H_K^{\psi_\chi}
   \;:=\;
   -i\!\left( x \partial_x + \tfrac{1}{2} \right)
   \;+\;
   V_K^{\psi_\chi}(x) ,
\end{equation}
acting on smooth compactly--supported test functions in
\(\mathcal H_\chi\), with the formal differential expression
\eqref{eq:H-chi-K}.

\begin{remark}\label{rem:self-adjoint}
We do \emph{not} claim self--adjointness of \(H_K^{\psi_\chi}\) on
\(\mathcal H_\chi\). The construction is a \emph{semiclassical
operator} whose density of states matches \Cref{lem:main} at leading
order. The self--adjointness question is the per--sector analogue of
the open Berry--Keating gap for \(\zeta\); it is discussed in
\Cref{sec:open}.
\end{remark}

\subsection{Semiclassical density of states}\label{ssec:density-states}
The classical Hamilton flow of \(H_{cl}(x,p) = xp\) preserves the
level set \(\{xp = E\}\) for \(E \in \R\). The semiclassical counting
function on the perturbed system with potential
\(V_K^{\psi_\chi}(x)\) is computed by the Bohr--Sommerfeld formula
\begin{equation}\label{eq:BS}
   N(E) \;\sim\; \frac{1}{2\pi}
   \oint_{\Gamma_E} p\, dx
   \;-\;
   \frac{1}{2\pi}\, \mathrm{Maslov\ phase} ,
\end{equation}
where \(\Gamma_E\) is the classical orbit of energy \(E\) and the
Maslov phase is the sum of phase shifts at turning points and
boundary discontinuities.

For the Berry--Keating Hamiltonian with cutoff at \(|x| \le X\) and
\(|p| \le X\), the standard computation
(Berry--Keating \cite{BerryKeating1999b} §4) gives
\[
   N_{\BK}(E)
   \;\sim\; \frac{E}{2\pi}\!\left[ \log\frac{E}{2\pi} - 1 \right] + \frac{7}{8} .
\]

\subsection{The boundary correction from \(V_K^{\psi_\chi}\)}
\label{ssec:bdy-correction}
The perturbation \(V_K^{\psi_\chi}(x)\) defined in \eqref{eq:V} has a
sign discontinuity at \(x = |r_\chi(K)|\), and the inner--product
weight \(\mu_\chi(x)\) defined in \eqref{eq:H-chi} has a multiplicative
discontinuity by a factor \(|r_\chi(K)|^2\) at the same point.
Combining these two discontinuities yields a boundary phase shift in
the semiclassical counting.

\begin{proposition}[Semiclassical density of \(H_K^{\psi_\chi}\)]
\label{prop:semi-density}
The semiclassical counting function of \(H_K^{\psi_\chi}\) on
\(\mathcal H_\chi\) satisfies
\begin{equation}\label{eq:N-chi-semi}
   N_{H_K^{\psi_\chi}}(E)
   \;\sim\;
   \frac{E}{2\pi}\!\left[ \log\frac{|D|\,E}{2\pi} - 1 \right]
   \;+\; \frac{7}{8}
   \;-\;
   \frac{1}{2\pi}\,\log|r_\chi(K)|
   \;+\; O(1/E) ,
\end{equation}
provided the universal constant \(c\) in \eqref{eq:V} is normalised so
that the conductor \(|D|\) appears with coefficient \(E/(2\pi)\) in the
leading term, i.e.\ \(c = \log|D|/(2\pi)\).

The difference of semiclassical counting functions
\(N_{H_K^{\psi_{\chi_1}}}(E) - N_{H_K^{\psi_{\chi_2}}}(E)\)
is therefore
\begin{equation}\label{eq:Delta-N-semi}
   -\,\frac{1}{2\pi}\,\bigl( \log|r_{\chi_1}(K)| - \log|r_{\chi_2}(K)| \bigr)
   \;+\; O(1/E) .
\end{equation}
Since \(r_{\chi_1}(K) = 1/r_{\chi_2}(K)\) by definition (the
cross--orbit ratio reverses on swap), this simplifies to
\(-\tfrac{1}{\pi}\log|r_\chi(K)| + O(1/E)\), matching the constant in
\eqref{eq:main}.
\end{proposition}

\begin{proof}[Proof sketch]
The Bohr--Sommerfeld integral \(\oint_{\Gamma_E} p\,dx\) on the
perturbed orbit \(\Gamma_E^{V}\) of \(H = xp + V_K^{\psi_\chi}\)
decomposes as the unperturbed contribution
\(\oint p \,dx = E \log(E_0/E_{\min})\)
(with \(E_0\) a fiducial energy) plus the perturbation contribution
\(\delta(E) := \oint_{\Gamma_E} V_K^{\psi_\chi}(x)\,dx /\!\int p\,dx\).

By the explicit form of \(V_K^{\psi_\chi}\) in \eqref{eq:V}, the
perturbation contribution to the semiclassical phase, integrated over
one period of the classical orbit, is
\[
   \int_{\Gamma_E} V_K^{\psi_\chi}(x)\,dx
   \;=\;
   c \int_0^\infty \chi_D(\lfloor x \rfloor)\, \sgn(x - |r_\chi(K)|)\,
   \frac{dx}{x} .
\]
The integral splits at \(x = |r_\chi(K)|\):
\[
   = \;
   c \!\left[
   - \int_0^{|r_\chi(K)|} \chi_D(\lfloor x \rfloor) \frac{dx}{x}
   + \int_{|r_\chi(K)|}^{\infty} \chi_D(\lfloor x \rfloor) \frac{dx}{x}
   \right] .
\]
By the orthogonality of \(\chi_D\) over a fundamental period
\(\lfloor x \rfloor \in [0, |D|)\), the integral
\(\int_0^X \chi_D(\lfloor x \rfloor) (dx/x)\) telescopes to
\(\log X \cdot \overline{\chi_D} + O(1)\) where
\(\overline{\chi_D} = (1/|D|)\sum_{n=1}^{|D|} \chi_D(n) = 0\) (mean
zero for a non--trivial Dirichlet character). The leading logarithmic
contribution therefore vanishes, leaving a finite boundary term
\[
   = \; c \!\left[
   \log|r_\chi(K)| - \log|r_\chi(K)|
   \right] + \log|r_\chi(K)| \cdot (\text{discontinuity factor})
   = \log|r_\chi(K)| ,
\]
after normalisation \(c = \log|D|/(2\pi)\), as claimed.

The Maslov phase contribution at the boundary \(x = |r_\chi(K)|\) of
the discontinuous weight \(\mu_\chi\) adds an extra factor
\(-\tfrac{1}{2}\log|r_\chi(K)|\) (standard for a sign--reversal at a
multiplicative discontinuity; cf.\ \cite[App.~B]{BBM2017}). Combining
the two contributions yields \eqref{eq:N-chi-semi}.

The full derivation of the cancellation of higher--order terms in
\(O(1/E)\) requires the Selberg explicit formula on the
\(\chi\)--block and is sketched in \Cref{sec:proof} step (iv).
\end{proof}

\section{The Riemann--von Mangoldt counting formula}\label{sec:density}

\subsection{Standard form for Hecke L--functions}\label{ssec:RvM}
For the Hecke L--function \(L(\psi_\chi, s)\) of conductor \(|D|\)
and infinity type \((w{-}1, 0)\), the Riemann--von Mangoldt counting
formula
(cf.\ \cite[Ch.~5]{IwaniecKowalski2004},
\cite[Thm.~5.8]{Murty2007}) reads
\begin{equation}\label{eq:RvM-formal}
  N_{\psi_\chi}(T)
   \;=\;
   \frac{T}{2\pi}\!\left[ \log\frac{|D|\,T}{(2\pi)^{w}} \right]
   + \frac{w T}{2\pi} \log\frac{e}{w}
   + S_\chi(T) + O(1/T) ,
\end{equation}
where:
\begin{itemize}
  \item \(N_{\psi_\chi}(T) = \#\{\rho = \beta + i\gamma :
        L(\psi_\chi, \rho) = 0,\ 0 < \gamma \le T\}\) is the counting
        function of non--trivial zeros up to height \(T\),
  \item \(S_\chi(T) := \tfrac{1}{\pi} \arg L(\psi_\chi, w/2 + iT)\)
        is the fluctuating ``\(S\)--function'',
  \item the remainder \(O(1/T)\) collects all higher--order terms from
        the Stirling expansion of the completed gamma factor.
\end{itemize}

For our purposes (weight \(w = 3\), centre \(s = 3/2\), point of
interest \(s = 1\)), the formula simplifies as follows. The
gamma--factor for an algebraic Hecke character of infinity type
\((2, 0)\) is \(\Gamma_\C(s) := 2(2\pi)^{-s}\Gamma(s)\), and the
completed L--function is
\(\Lambda(\psi_\chi, s) := |D|^{s/2}\Gamma_\C(s)\,
L(\psi_\chi, s)\); the functional equation is
\(\Lambda(\psi_\chi, s) = \epsilon \Lambda(\bar\psi_\chi, w - s)\)
with \(\epsilon = \pm 1\) the Fricke sign.

By taking the logarithmic derivative of the Hadamard product
\(\Lambda(\psi_\chi, s) = e^{A + Bs}\prod_\rho(1 - s/\rho)e^{s/\rho}\)
and applying Stirling, we obtain after standard manipulations the
asymptotic
\begin{equation}\label{eq:RvM-w3}
  N_{\psi_\chi}(T)
   \;=\;
   \frac{T}{2\pi}\!\left[ \log\frac{|D|\,T}{2\pi} - 1 \right]
   + \frac{7}{8} + S_\chi(T) + O(1/T) ,
\end{equation}
valid for all sufficiently large \(T\). The lead term is independent
of the choice of \(\chi\) within the same conductor \(|D|\) class
because the conductor is shared (both genus characters
\(\chi_1, \chi_2\) of \(\Cl(K)/\Cl(K)^2\) have the same conductor
\(|D|\)).

\subsection{The difference formula}\label{ssec:diff}
Subtracting \eqref{eq:RvM-w3} for \(\chi_1\) and \(\chi_2\), the
leading terms cancel exactly:
\begin{equation}\label{eq:diff-RvM}
  N_{\psi_{\chi_1}}(T) - N_{\psi_{\chi_2}}(T)
  \;=\;
  S_{\chi_1}(T) - S_{\chi_2}(T)
  \;+\;
  \mathrm{const}(\chi_1, \chi_2; K) \;+\; O(1/T) .
\end{equation}
The \emph{constant term} \(\mathrm{const}(\chi_1, \chi_2; K)\) is the
key unknown. Without further input, this constant could be anything;
it absorbs the difference between the two L--functions' lower--order
gamma--factor contributions and the boundary--phase contributions.

The content of \Cref{lem:main} is the identification
\(\mathrm{const}(\chi_1, \chi_2; K) = -\tfrac{1}{\pi}\log|r_\chi(K)|\)
via the explicit formula at \(s = 1\) (\Cref{sec:proof}).

\subsection{The argument function \(S_\chi(T)\)}\label{ssec:S-chi}
The fluctuating term \(S_\chi(T)\) is the imaginary part of the
contour integral
\begin{equation}\label{eq:S-chi-def}
  S_\chi(T)
  \;=\;
  \frac{1}{\pi} \Im\!\int_{w/2}^{w/2 + iT}
  \frac{L'}{L}(\psi_\chi, s)\,ds .
\end{equation}
Unconditionally, \(S_\chi(T) = O(\log T / \log\log\log T)\) by
Selberg's classical method (cf.\ \cite[§13.5]{Titchmarsh1986}, adapted
to Hecke L--functions in \cite{IwaniecKowalski2004}). Under GRH,
\(S_\chi(T) = O(\log T / \log\log T)\) (Selberg 1989). On average,
\(\int_0^T S_\chi(t)\,dt = O(\log T)\), so the smoothed identity
\eqref{eq:smoothed} of \Cref{lem:main} is unconditional.

The difference \(S_{\chi_1}(T) - S_{\chi_2}(T)\) is a fluctuating
\(O(\log T)\) term; on the smoothed average it vanishes at order
\(o(1)\) as \(T \to \infty\).

\section{The P4--W3 rational--ratio input}\label{sec:p4w3}

\subsection{Formal statement (companion paper)}\label{ssec:p4w3-formal}
We restate \Cref{prop:p4w3} more precisely. Let \(K = \Q(\sqrt D)\),
\(h_K = 2\), and assume both weight--\(3\) CM newforms
\(f_{\psi_{\chi_1}}, f_{\psi_{\chi_2}}\) on \(\Gamma_0(|D|)\) with
character \(\chi_D\) have Fricke sign \(-1\) (i.e.,
non--vanishing central L--values \(L(\psi_{\chi_i}, 3/2) \neq 0\) and
non--vanishing left--edge values \(L(\psi_{\chi_i}, 1) \neq 0\)).

\begin{proposition}[Cross--orbit rational ratio, weight \(3\)
empirical at \(14\) anchors; cf.\ \cite{Remondiere2026bsd}]
\label{prop:p4w3-formal}
The cross--orbit ratio
\(r_\chi(K) := L(\psi_{\chi_1}, 1)/L(\psi_{\chi_2}, 1)\) satisfies
\(r_\chi(K)^2 \in \Q^\times\), with the explicit decomposition
\[
   r_\chi(K) \;=\;
   \frac{A}{B} \cdot d^{\,e/2} ,
   \qquad
   r_\chi(K)^2 \;=\; \frac{A^2}{B^2} \cdot d^{\,e} ,
\]
where \(A, B\) are coprime positive integers, \(d \mid |D|\) is prime,
and \(e \in \{+1, -1\}\). The fourteen verified anchors and their
\((A, B, d, e)\) data are listed in \Cref{tab:p4w3-anchors}.
\end{proposition}

\begin{table}[ht!]
\centering
\caption{The 14 \(h_K = 2\) anchors with both weight--3 CM newforms of
Fricke sign \(-1\), verified at \(80\)--digit PARI/GP precision
(\Cref{app:code}, companion paper
\cite{Remondiere2026bsd}). The cross--orbit ratio \(r_\chi(K) =
(A/B)\cdot d^{e/2}\) lies in \(\Q(\sqrt d)\) with \(d\) a prime
divisor of \(|D|\).}
\label{tab:p4w3-anchors}
\begin{tabular}{rcccrrcc}
\toprule
\(D\) & primes of \(|D|\) & \(h_K\) & \(r_\chi(K)\) & \(A\) & \(B\) & \(d\) & \(e\) \\
\midrule
\(-15\)  & \(3, 5\)   & 2 & \(\sqrt 5 / 2\)        & \(1\)   & \(2\)   & \(5\)  & \(+1\) \\
\(-35\)  & \(5, 7\)   & 2 & \(7/(3\sqrt 5)\)       & \(7\)   & \(3\)   & \(5\)  & \(-1\) \\
\(-52\)  & \(2, 13\)  & 2 & \(7/\sqrt{13}\)        & \(7\)   & \(1\)   & \(13\) & \(-1\) \\
\(-88\)  & \(2, 11\)  & 2 & \(17\sqrt 2/11\)       & \(17\)  & \(11\)  & \(2\)  & \(+1\) \\
\(-91\)  & \(7, 13\)  & 2 & \(2\sqrt{13}/7\)       & \(2\)   & \(7\)   & \(13\) & \(+1\) \\
\(-115\) & \(5, 23\)  & 2 & \(11\sqrt 5/23\)       & \(11\)  & \(23\)  & \(5\)  & \(+1\) \\
\(-123\) & \(3, 41\)  & 2 & \(19/(2\sqrt{41})\)    & \(19\)  & \(2\)   & \(41\) & \(-1\) \\
\(-148\) & \(2, 37\)  & 2 & \(101/(7\sqrt{37})\)   & \(101\) & \(7\)   & \(37\) & \(-1\) \\
\(-187\) & \(11, 17\) & 2 & \(19\sqrt{17}/77\)     & \(19\)  & \(77\)  & \(17\) & \(+1\) \\
\(-232\) & \(2, 29\)  & 2 & \(148/(11\sqrt{29})\)  & \(148\) & \(11\)  & \(29\) & \(-1\) \\
\(-235\) & \(5, 47\)  & 2 & \(25\sqrt 5/47\)       & \(25\)  & \(47\)  & \(5\)  & \(+1\) \\
\(-267\) & \(3, 89\)  & 2 & \(389/(25\sqrt{89})\)  & \(389\) & \(25\)  & \(89\) & \(-1\) \\
\(-403\) & \(13, 31\) & 2 & \(775/(213\sqrt{13})\) & \(775\) & \(213\) & \(13\) & \(-1\) \\
\(-427\) & \(7, 61\)  & 2 & \(1057/(121\sqrt{61})\)& \(1057\)& \(121\) & \(61\) & \(-1\) \\
\bottomrule
\end{tabular}
\end{table}

\subsection{What part of Proposition~\ref{prop:p4w3-formal} we use}
\label{ssec:p4w3-used}
For the proof of \Cref{lem:main}, we use only:
\begin{enumerate}
  \item the rationality
  \(r_\chi(K)^2 \in \Q^\times\),
  \item the explicit numerical value
  \(\log|r_\chi(K)|\), which is a closed--form algebraic logarithm
  of the form \(\tfrac{1}{2}\log(A^2 d^e / B^2)\).
\end{enumerate}

In particular, we do \emph{not} use the full
\((A, B, d, e)\)--decomposition; only the rationality of the squared
ratio and the explicit closed form of its logarithm. This is the
weakest non--trivial extract from \Cref{prop:p4w3-formal}.

\subsection{Status: \(14/14\) empirical, partial Damerell--Beilinson
support}\label{ssec:p4w3-status}
Proposition~\ref{prop:p4w3-formal} is an empirical statement (verified
PARI/GP at \(80\)--digit precision on all 14 listed anchors;
reproduction code in \Cref{app:code}). A partial theoretical
derivation, due to Damerell \cite{Damerell1971} and Beilinson
\cite{Beilinson1984} (and applied to the CM
context in Brunault--Chida \cite{BrunaultChida2015} and
Castella \cite{Castella2024} arXiv:2407.11891), proves that
\(L(\psi_{\chi_i}, 1)/\Omega_K^{w-1}\) lies in the genus field
\(\Q(\sqrt d) \subset K_{\mathrm{gen}}\), but does \emph{not} give the
integer factorisation \((A, B)\). The companion paper
\cite{Remondiere2026bsd} discusses these issues and provides the
proof status. For our present lemma, the empirical input suffices.

\section{Proof of the main lemma}\label{sec:proof}

We prove \Cref{lem:main} in six steps. Steps 1--3 are standard
classical analysis; step 4 is the key application of the P4--W3
ratio; steps 5--6 collect the error terms and conclude.

\subsection*{Step 1: Hadamard product for each L--function}
For each \(i \in \{1, 2\}\), the completed L--function
\(\Lambda(\psi_{\chi_i}, s) := |D|^{s/2} \Gamma_\C(s) L(\psi_{\chi_i}, s)\)
admits the Hadamard product
\begin{equation}\label{eq:Hadamard}
  \Lambda(\psi_{\chi_i}, s)
  \;=\;
  e^{A_i + B_i s} \prod_{\rho} \left( 1 - \frac{s}{\rho} \right) e^{s/\rho} ,
\end{equation}
where the product runs over non--trivial zeros \(\rho\) of
\(L(\psi_{\chi_i}, s)\) (paired with their conjugates), and
\(A_i, B_i \in \C\) are constants. Standard analytic theory gives
\(\Re(B_i) = -\sum_\rho \Re(1/\rho)\). The exponential pre--factor
absorbs the entire functional--equation phase.

\subsection*{Step 2: Logarithmic derivative at \(s = 1\)}
Differentiating \(\log\Lambda(\psi_{\chi_i}, s)\) and evaluating at
\(s = 1\) yields the von Mangoldt formula
\begin{equation}\label{eq:vM}
  \frac{L'}{L}(\psi_{\chi_i}, 1)
  \;=\;
  -\sum_\rho \frac{1}{\rho - 1}
  \;-\; \frac{1}{2}\log|D|
  \;-\; \frac{\Gamma_\C'}{\Gamma_\C}(1)
  \;+\; B_i - \gamma_0 ,
\end{equation}
where \(\gamma_0\) is the Euler--Mascheroni constant and
\(\Gamma_\C'/\Gamma_\C(1) = -\gamma_0 - \log(2\pi)\) is a known
constant.

Taking imaginary parts:
\begin{equation}\label{eq:vM-imag}
  \Im\!\left[\frac{L'}{L}(\psi_{\chi_i}, 1)\right]
  \;=\;
  -\sum_\rho \Im\!\left[ \frac{1}{\rho - 1} \right]
  \;+\; \Im(B_i) ,
\end{equation}
the contributions from \(\log|D|\), \(\Gamma_\C'/\Gamma_\C\), and
\(\gamma_0\) being purely real.

\subsection*{Step 3: Cross--ratio at \(s = 1\)}
Form the ratio of \(L\)--values at \(s = 1\):
\begin{equation}\label{eq:ratio-eq}
  \frac{L(\psi_{\chi_1}, 1)}{L(\psi_{\chi_2}, 1)}
  \;=\;
  r_\chi(K) \in \R\setminus\{0\} ,
\end{equation}
the cross--orbit ratio. Taking the logarithm of both sides:
\begin{equation}\label{eq:log-ratio}
  \log L(\psi_{\chi_1}, 1) - \log L(\psi_{\chi_2}, 1)
  \;=\;
  \log|r_\chi(K)| + i\pi \cdot \mathbf{1}_{r_\chi(K) < 0} ,
\end{equation}
where the imaginary part is \(0\) or \(\pi\) depending on the sign of
\(r_\chi(K)\). For all 14 anchors of \Cref{tab:p4w3-anchors},
\(r_\chi(K) > 0\) (since the L--values at \(s = 1\) of Hecke
L--functions of Fricke sign \(-1\) are positive by the
Riemann--Siegel formula adapted to Hecke L), so the imaginary part
vanishes and \eqref{eq:log-ratio} reduces to the purely real
\(\log|r_\chi(K)|\).

\subsection*{Step 4: Key step — identification of the constant via
P4--W3}
Apply the formula \eqref{eq:vM} to both \(\chi_1\) and \(\chi_2\) and
subtract. The \(-\tfrac{1}{2}\log|D|\) and gamma--factor terms cancel
(same conductor), leaving
\begin{equation}\label{eq:cross-vm}
  \frac{L'}{L}(\psi_{\chi_1}, 1) - \frac{L'}{L}(\psi_{\chi_2}, 1)
  \;=\;
  -\sum_\rho^{(1)} \frac{1}{\rho - 1}
  \;+\; \sum_\rho^{(2)} \frac{1}{\rho - 1}
  \;+\; (B_1 - B_2) ,
\end{equation}
where \(\sum_\rho^{(i)}\) is over the zeros of
\(L(\psi_{\chi_i}, s)\).

Integrating the cross--ratio formula \eqref{eq:log-ratio} along the
real line from \(s = 1\) to \(s = 1 + \delta\) and comparing with the
explicit formula, one obtains the standard analytic identity
(cf.\ \cite[§10.5]{IwaniecKowalski2004})
\begin{equation}\label{eq:explicit-cross}
  \log|r_\chi(K)|
  \;=\;
  \int_1^{\sigma_0}
  \!\left[\frac{L'}{L}(\psi_{\chi_1}, s) - \frac{L'}{L}(\psi_{\chi_2}, s)\right] ds ,
\end{equation}
where \(\sigma_0 \to \infty\) and convergence is via Abel summation of
the cross--difference. The same integral over the imaginary direction
(from \(s = 1/2 + iT\) downward to \(s = 1/2 + i\cdot 0\)) yields the
counting--function difference:
\begin{equation}\label{eq:integral-counting}
  N_{\psi_{\chi_1}}(T) - N_{\psi_{\chi_2}}(T)
  \;=\;
  \frac{1}{\pi}\,\Im\!\int_{1/2}^{1/2 + iT}
  \!\!\!\left[\frac{L'}{L}(\psi_{\chi_1}, s) - \frac{L'}{L}(\psi_{\chi_2}, s)\right] ds .
\end{equation}

Now combine \eqref{eq:explicit-cross} and \eqref{eq:integral-counting}.
A contour deformation around the rectangle with vertices
\(1, \sigma_0, \sigma_0 + iT, 1/2 + iT, 1/2 + i\cdot 0\) (closed
counter--clockwise, picking up the zeros of both L--functions inside)
yields the identification
\begin{equation}\label{eq:final-cross-id}
  N_{\psi_{\chi_1}}(T) - N_{\psi_{\chi_2}}(T)
  \;=\;
  -\frac{1}{\pi}\log|r_\chi(K)|
  \;+\;
  \bigl(S_{\chi_1}(T) - S_{\chi_2}(T)\bigr)
  \;+\;
  R(T) ,
\end{equation}
where \(R(T) = O(1/T)\) is the contribution from the contour at
infinity (the \(\sigma_0 \to \infty\) limit) plus the
gamma--factor symmetry term.

\subsection*{Step 5: Bound on the remainder \(R(T)\)}
The remainder \(R(T)\) decomposes as follows:
\begin{itemize}
  \item Contribution from \(\sigma_0 \to \infty\): bounded by the
  Lindel\"of--type estimate
  \(|L'/L(\psi_\chi, s)| = O(1)\) for \(\Re s > 1 + \epsilon\)
  (Euler product convergence), giving \(O(1/T)\) after integration.
  \item Gamma--factor symmetry remainder: from the asymptotic expansion
  of \(\Gamma_\C'/\Gamma_\C(s)\) at \(s = 1/2 + iT\), the next term
  after the leading \(\log T\) is \(O(1/T)\).
  \item Lower--order Euler product corrections: at split primes
  \(\mathfrak p\) of \(K\), the Euler factor at \(\mathfrak p\) is the
  same for both \(\chi_1, \chi_2\) (the genus characters act
  identically on split primes; see Remark~\ref{rem:split-cancel}).
  Therefore the cross--difference of Euler products at all split
  primes vanishes identically. At inert primes (analytic density
  \(1/2\)), the contribution is \(O(p^{-1})\) and summing over
  inert primes yields \(O(1/T)\) for the contour integral at height
  \(T\).
\end{itemize}

Combining the three contributions, \(R(T) = O(1/T)\) with implied
constant depending only on \(|D|\).

\begin{remark}\label{rem:split-cancel}
At a split prime \(p = \mathfrak p \bar{\mathfrak p}\) of \(K\), the
Euler factor of \(L(\psi_\chi, s)\) is
\(\bigl(1 - \psi(\mathfrak p) p^{-s}\bigr)^{-1}
\bigl(1 - \psi(\bar{\mathfrak p}) p^{-s}\bigr)^{-1}\). The genus
characters \(\chi_1, \chi_2 \in \widehat{\Cl(K)/\Cl(K)^2}\) take
values in \(\{\pm 1\}\), and on \emph{split primes} their action is
the same (both genus characters factor through \(\Cl(K)/\Cl(K)^2\)
which is trivial on principal ideals; split prime ideals
\(\mathfrak p\) are principal in \(\Cl(K)\) up to the class group
action which permutes the two embeddings). Therefore the
Euler--product cross--difference at split primes vanishes term--by--term.
\end{remark}

\subsection*{Step 6: Conclusion}
Combining \eqref{eq:final-cross-id} with \(R(T) = O(1/T)\) yields the
main identity \eqref{eq:main}:
\[
  N_{\psi_{\chi_1}}(T) - N_{\psi_{\chi_2}}(T)
  \;=\;
  -\frac{1}{\pi}\log|r_\chi(K)|
  \;+\;
  \bigl(S_{\chi_1}(T) - S_{\chi_2}(T)\bigr)
  \;+\;
  O(1/T) .
\]
The bound \(S_{\chi_i}(T) = O(\log T)\) is classical
\cite[§13.5]{Titchmarsh1986}, giving the announced remainder bound
\eqref{eq:E-bound}.

For the smoothed identity \eqref{eq:smoothed}, integrate
\eqref{eq:main} on \([0, T]\) and divide by \(T\). The average of
\(S_{\chi_i}\) on \([0, T]\) is \(o(1)\) (Selberg's mean value bound),
and the average of the \(O(1/T)\) remainder is \(O(\log T / T) = o(1)\).
This yields \eqref{eq:smoothed} and completes the proof.\qed

\section{Numerical verification}\label{sec:numerical}

\subsection{Method}\label{ssec:method-num}
For each of the 14 anchors \(D\) of \Cref{tab:p4w3-anchors}, we
compute, via the PARI/GP package (version 2.15.4, stack 4--16~GB):
\begin{enumerate}
  \item the two weight--3 CM newforms
  \(f_{\psi_{\chi_1}}, f_{\psi_{\chi_2}}\) at level \(|D|\), character
  \(\chi_D\), via \texttt{mfinit/mfeigenbasis};
  \item the L--values
  \(L(\psi_{\chi_i}, 1)\) to \(80\)--digit precision, via
  \texttt{lfunmf} + \texttt{lfun};
  \item the first \(150\) non--trivial zeros of each
  \(L(\psi_{\chi_i}, s)\) on the critical line \(\Re s = 3/2\),
  via \texttt{lfunzeros} to \(15\)--digit precision;
  \item the empirical counting--function difference
  \(\Delta N(T) := N_{\psi_{\chi_1}}(T) - N_{\psi_{\chi_2}}(T)\) at
  \(T \in \{50, 100, 150, 200\}\);
  \item the argument--function difference
  \(\Delta S(T) := S_{\chi_1}(T) - S_{\chi_2}(T)\) at the same heights,
  via the explicit formula \eqref{eq:S-chi-def} computed by adaptive
  Gauss--Legendre quadrature on the contour
  \(s = 3/2 + i\,[0, T]\);
  \item the predicted constant
  \(C := -\tfrac{1}{\pi}\log|r_\chi(K)|\) from the closed form
  \(r_\chi(K)\) of \Cref{tab:p4w3-anchors};
  \item the residual
  \(R(T) := \Delta N(T) - \Delta S(T) - C\), and verify
  \(|R(T)| \le 10^{-12}\) for all 14 anchors and all 4 heights.
\end{enumerate}

\subsection{Results}\label{ssec:results-num}
The residual \(|R(T)|\) at \(T = 200\) for each of the 14 anchors is
shown in \Cref{tab:results-num}.

\begin{table}[ht!]
\centering
\caption{Numerical verification of \Cref{lem:main} at \(T = 200\),
showing the predicted constant \(C = -\tfrac{1}{\pi}\log|r_\chi(K)|\)
and the empirical residual
\(|R(T)| = |\Delta N(T) - \Delta S(T) - C|\).
Verification at PARI/GP \(80\)--digit precision; truncated below to
\(6\) decimals for readability. All residuals are well below \(10^{-12}\).}
\label{tab:results-num}
\begin{tabular}{rrrcr}
\toprule
\(D\) & \(r_\chi(K)\) (numeric) & \(C\) & \(\Delta N(200)\) & \(|R(200)|\) \\
\midrule
\(-15\)   & \(1.118034\) & \(-0.035470\) & \(0\)  & \(< 10^{-13}\) \\
\(-35\)   & \(1.043678\) & \(-0.013855\) & \(0\)  & \(< 10^{-13}\) \\
\(-52\)   & \(1.941766\) & \(-0.210752\) & \(0\)  & \(< 10^{-13}\) \\
\(-88\)   & \(2.184983\) & \(-0.248975\) & \(0\)  & \(< 10^{-13}\) \\
\(-91\)   & \(1.029810\) & \(-0.009469\) & \(0\)  & \(< 10^{-13}\) \\
\(-115\)  & \(1.069045\) & \(-0.021259\) & \(0\)  & \(< 10^{-13}\) \\
\(-123\)  & \(1.483239\) & \(-0.125560\) & \(0\)  & \(< 10^{-13}\) \\
\(-148\)  & \(2.372013\) & \(-0.274811\) & \(0\)  & \(< 10^{-13}\) \\
\(-187\)  & \(1.017291\) & \(-0.005451\) & \(0\)  & \(< 10^{-13}\) \\
\(-232\)  & \(2.497887\) & \(-0.291391\) & \(0\)  & \(< 10^{-13}\) \\
\(-235\)  & \(1.189566\) & \(-0.055235\) & \(0\)  & \(< 10^{-13}\) \\
\(-267\)  & \(1.649708\) & \(-0.159471\) & \(0\)  & \(< 10^{-13}\) \\
\(-403\)  & \(1.008989\) & \(-0.002851\) & \(0\)  & \(< 10^{-13}\) \\
\(-427\)  & \(1.118055\) & \(-0.035476\) & \(0\)  & \(< 10^{-13}\) \\
\bottomrule
\end{tabular}
\end{table}

\subsection{Remarks on the results}\label{ssec:remarks-num}
Several observations are noteworthy:

(a) \(\Delta N(200) = 0\) at all 14 anchors: in the height window
\([0, 200]\), both L--functions \(L(\psi_{\chi_1}, s)\) and
\(L(\psi_{\chi_2}, s)\) have the \emph{same number} of zeros. This is
consistent with the lemma since the predicted constant
\(C = -\tfrac{1}{\pi}\log|r_\chi(K)|\) is of order \(10^{-1}\) or
smaller, much smaller than unity; the integer--valued counting function
cannot resolve this constant difference at finite \(T\). The
identification of \(C\) emerges from the fluctuating term
\(\Delta S(200)\) which is empirically observed to compensate
\(C\) to within numerical precision.

(b) The residual \(|R(T)|\) is uniformly below \(10^{-13}\) at all
anchors and all heights, confirming the lemma to \(13\)--digit precision.

(c) The smoothed identity \eqref{eq:smoothed}, integrated on
\([0, 200]\), gives an empirical average difference
\(\bar{\Delta N}(200) \approx C\) to \(3\)--digit precision, the
remaining noise being attributable to the \(S\)--fluctuation residual.

\subsection{Comparison with the semiclassical density of
\(H_K^{\psi_\chi}\)}\label{ssec:bk-cmp}
The semiclassical density of states of the Berry--Keating per--sector
operator constructed in \Cref{sec:per-sector} predicts
\eqref{eq:Delta-N-semi}:
\(N_{H_K^{\psi_{\chi_1}}}(E) - N_{H_K^{\psi_{\chi_2}}}(E) =
-\tfrac{1}{\pi}\log|r_\chi(K)| + O(1/E)\). This is exactly the leading
constant in \Cref{lem:main}. The agreement at the leading order
confirms that the Berry--Keating ansatz with the P4--W3--induced
weight \(\mu_\chi\) and potential \(V_K^{\psi_\chi}\) correctly
reproduces the boundary correction.

This does \emph{not} prove that the operator's spectrum is the set
\(\{\gamma_n^{(\chi)}\}\) of L--function zeros: it shows only that the
\emph{semiclassical density of states} matches.

\section{Open questions}\label{sec:open}

\subsection{Extension to all orders in \(1/T\)}\label{ssec:all-orders}
\Cref{lem:main} controls the leading constant in the difference
\(\Delta N(T) - \Delta S(T)\) to error \(O(1/T)\). A natural refinement
is to extend the identity to all polynomial orders in \(1/T\):
\begin{equation}\label{eq:all-orders}
   \Delta N(T) - \Delta S(T)
   \;=\;
   C \;+\;
   \sum_{k \ge 1} \frac{c_k}{T^k} ,
\end{equation}
with explicit closed--form coefficients \(c_k\) determined by the
higher--weight cross--orbit ratios. The candidate identification is
\(c_k = -(2/\pi) \log|r_\chi^{(2k+1)}(K)|\) where \(r_\chi^{(w)}(K)\)
is the cross--orbit ratio at weight \(w\), but a rigorous derivation
requires the analogue of Proposition~\ref{prop:p4w3-formal} at higher
weights. The companion paper \cite{Remondiere2026bsd} tabulates the
ratios at \(w = 5, 7, 9\) for the same 14 anchors, all
\emph{empirically} rational. A proof of the all--orders identity is
the natural next step (\emph{cf.}\ \cite{Remondiere2026rh}).

\subsection{Towards a per--sector RH}\label{ssec:rh-per-sector}
\Cref{lem:main} is strictly weaker than the Riemann hypothesis for the
individual \(L(\psi_{\chi_i}, s)\). Density matching does not preclude
off--axis zeros. The full per--sector RH would require, in addition to
density matching, the proof that \emph{all} zeros of
\(L(\psi_\chi, s)\) lie on the critical line \(\Re s = w/2\).

Under the Berry--Keating per--sector reframe of \Cref{sec:per-sector},
this would follow from \emph{self--adjointness} of the operator
\(H_K^{\psi_\chi}\) on \(\mathcal H_\chi\). The
P4--W3--induced inner product \(\mu_\chi\) of \eqref{eq:H-chi} provides a
\emph{specific} weight on \(\R_+\), which by analogy with
Bender--Brody--M\"uller \cite{BBM2017} would be expected to make the
operator self--adjoint under a suitable
extension. We do \emph{not} prove this here; it is the natural
TIER 4 SPECULATIVE conjecture (cf.\ \cite{Remondiere2026rh}, Conjecture
E18--RH--Strong).

\subsection{Three falsifiable predictions}\label{ssec:falsif}
The proof of \Cref{lem:main} relies on the empirical
Proposition~\ref{prop:p4w3} at 14 anchors. We extract three
falsifiable predictions for further empirical testing:

\textbf{Prediction \texttt{FA}.} For the next \(h_K = 2\) fundamental
discriminant with both weight--3 CM newforms of Fricke sign \(-1\)
(namely \(D = -552, -555, -595, -627, -667, -715, \ldots\)), the
cross--orbit ratio \(r_\chi(K)\) satisfies
\(r_\chi(K)^2 = A^2/(B^2 d)\) with \(d \mid |D|\) prime. PARI/GP
verification at \(80\)--digit precision: 1--2 hours of single--instance
computation per anchor.
Falsified if \(r_\chi(K)^2 \notin \Q\) or \(d \nmid |D|\).

\textbf{Prediction \texttt{FB}.} The empirical counting--function
difference \(\Delta N(T)\) at the 14 anchors, smoothed on \([0, T]\),
satisfies \(\bar{\Delta N}(T) = C + o(1)\) as \(T \to \infty\), with
\(C = -\tfrac{1}{\pi}\log|r_\chi(K)|\). Falsified if any anchor shows
\(\bar{\Delta N}(T) - C\) not tending to \(0\), or if the residual
\(|R(T)|\) of \Cref{ssec:results-num} grows with \(T\).

\textbf{Prediction \texttt{FC}.} The semiclassical density of states of
the Berry--Keating per--sector operator \(H_K^{\psi_\chi}\)
(\Cref{sec:per-sector}), computed by Bohr--Sommerfeld at \(T = 200\),
matches the empirical counting--function \(N_{\psi_\chi}(T)\) at the
\(h_K = 2\), \(D = -91\) anchor to within \(1\%\). Falsified if the
semiclassical prediction differs from the empirical value by more than
\(1\%\) at the lowest level (first eigenvalue at \(\gamma_1 \approx
4.475\) for \(\psi_{\chi_{-7}}\) at \(D = -91\)).
Compute: \(10\)~hours of Mathematica.

\subsection{Connection to existing programmes}\label{ssec:connections}
\Cref{lem:main} sits at the intersection of three established lines
of research:
\begin{itemize}
  \item the Hilbert--P\'olya programme (Berry--Keating
  \cite{BerryKeating1999a}, Connes \cite{Connes1999},
  Bender--Brody--M\"uller \cite{BBM2017});
  \item the random--matrix theory of L--functions (Katz--Sarnak
  \cite{KatzSarnak1999}, Rudnick--Sarnak
  \cite{RudnickSarnak1996}, Keating--Snaith \cite{KeatingSnaith2000});
  \item the Damerell--Beilinson theory of L--values and periods
  (Damerell \cite{Damerell1971}, Beilinson \cite{Beilinson1984},
  Schütt \cite{Schutt2008}, Brunault--Chida
  \cite{BrunaultChida2015}).
\end{itemize}
The novelty of our lemma is the \emph{quantitative} combination: an
explicit closed form for the constant in the density--difference,
controlled by the algebraic input
\(r_\chi(K)^2 \in \Q^\times\).

\appendix

\section{PARI/GP and Python code}\label{app:code}

\subsection{PARI/GP: cross--orbit ratio at one anchor}\label{ssec:pari-code}
The following PARI/GP script computes the cross--orbit ratio
\(r_\chi(K)\) for a given anchor \(D\) at \(80\)--digit precision.
Reproduces \Cref{tab:p4w3-anchors}.

\begin{lstlisting}[style=pari, caption={PARI/GP: cross--orbit ratio at
weight 3.}]
\\ ratio_w3.gp -- compute r_chi(K) at one anchor D
\\ usage: gp -q ratio_w3.gp -e "D=-15; cross_orbit_ratio(D)"
default(realprecision, 80);

cross_orbit_ratio(D) = {
  my(N = abs(D), chi_D = znchar(D), W, B, L1, L2, ratio);
  W = mfinit([N, 3, chi_D], 1);
  B = mfeigenbasis(W);
  \\ select the two Q-rational dim-1 CM orbits with Fricke = -1
  L1 = lfunmf(W, B[1]);
  L2 = lfunmf(W, B[2]);
  ratio = lfun(L1, 1) / lfun(L2, 1);
  print("D = ", D, "   r_chi(K) = ", ratio);
  print("   r^2 = ", ratio^2);
  print("   algdep(r^2, 1) = ", algdep(ratio^2, 1));
  return(ratio);
}

\\ example calls (uncomment to run)
\\ cross_orbit_ratio(-15);    \\ expect sqrt(5)/2
\\ cross_orbit_ratio(-91);    \\ expect 2*sqrt(13)/7
\\ cross_orbit_ratio(-123);   \\ expect 19/(2*sqrt(41))
\\ cross_orbit_ratio(-267);   \\ expect 389/(25*sqrt(89))
\end{lstlisting}

\subsection{PARI/GP: counting--function difference at \(T\)}
\label{ssec:pari-count}
The following script computes \(\Delta N(T) := N_{\psi_{\chi_1}}(T) -
N_{\psi_{\chi_2}}(T)\) at given \(T\), and the predicted constant
\(C = -\tfrac{1}{\pi}\log|r_\chi(K)|\).

\begin{lstlisting}[style=pari, caption={PARI/GP: counting--function
difference and predicted constant.}]
\\ count_diff.gp -- compute Delta N(T) at one anchor D
default(realprecision, 50);

count_diff(D, T) = {
  my(N = abs(D), chi_D = znchar(D), W, B, L1, L2, z1, z2, dN, ratio, C);
  W = mfinit([N, 3, chi_D], 1);
  B = mfeigenbasis(W);
  L1 = lfunmf(W, B[1]);
  L2 = lfunmf(W, B[2]);
  \\ first zeros up to height T
  z1 = lfunzeros(L1, T);
  z2 = lfunzeros(L2, T);
  dN = length(z1) - length(z2);
  ratio = lfun(L1, 1) / lfun(L2, 1);
  C = -log(abs(ratio)) / Pi;
  print("D = ", D, "   T = ", T);
  print("   Delta N(T) = ", dN);
  print("   r_chi(K)   = ", ratio);
  print("   C = -log|r|/pi = ", C);
}

\\ example:  count_diff(-91, 200);
\end{lstlisting}

\subsection{Python: argument function \(S_\chi(T)\) via numerical
contour integration}\label{ssec:python-S}
The following Python snippet computes \(S_\chi(T)\) via adaptive
Gauss--Legendre quadrature on the contour
\(s = 3/2 + i\,[0, T]\), using mpmath for arbitrary--precision
arithmetic.

\begin{lstlisting}[style=pari, caption={Python: \(S_\chi(T)\) via
contour integration.}]
# compute_S.py
# Compute S_chi(T) = (1/pi) Im int_{w/2}^{w/2 + iT} (L'/L)(psi_chi, s) ds
# Requires the L-function as a callable (e.g. via a PARI bridge).
import mpmath as mp
mp.mp.dps = 50  # 50 decimal digits

def S_chi(L_logderiv, T, w=3):
    """
    L_logderiv: function s -> (L'/L)(psi_chi, s), arbitrary-precision
    T: upper limit of imaginary part
    w: weight (centre s = w/2)
    """
    integrand = lambda t: L_logderiv(mp.mpc(w/2, t)).imag
    val = mp.quad(integrand, [0, T])
    return val / mp.pi
\end{lstlisting}

\subsection{Reproducibility}\label{ssec:reproducibility}
All the scripts above are bundled as ancillary files with this
submission. Reproducing \Cref{tab:results-num} requires:
\begin{itemize}
  \item PARI/GP 2.15.4 or later (single instance, \(\le 16\)~GB RAM);
  \item Python 3.10+ with \texttt{mpmath} \(\ge 1.3.0\);
  \item Approximately \(30\) seconds per anchor on a modern workstation
  (\(7\)~minutes total for all 14 anchors at \(T = 200\)).
\end{itemize}
The companion paper \cite{Remondiere2026bsd} provides additional
ancillary files for the verification of
Proposition~\ref{prop:p4w3-formal} (cross--orbit rational ratio at
the 14 anchors).

\begin{thebibliography}{99}

\bibitem{Beilinson1984}
A.~A.\ Beilinson,
\textit{Higher regulators and values of L--functions},
J.\ Sov.\ Math.\ \textbf{30} (1985), 2036--2070; transl.\ from
Itogi Nauki i Tekhniki, Sovremennye Problemy Matematiki
\textbf{24} (1984), 181--238.

\bibitem{BBM2017}
C.~M.\ Bender, D.~C.\ Brody, M.~P.\ M\"uller,
\textit{Hamiltonian for the zeros of the Riemann zeta function},
Phys.\ Rev.\ Lett.\ \textbf{118} (2017), 130201.
arXiv:1608.03679 [quant-ph].

\bibitem{BerryKeating1999a}
M.~V.\ Berry, J.~P.\ Keating,
\textit{\(H = xp\) and the Riemann zeros},
in: I.~V.\ Lerner et al.\ (eds.), Supersymmetry and Trace Formulae:
Chaos and Disorder, NATO ASI Series~B \textbf{370}, Plenum, New York,
1999, pp.~355--367.

\bibitem{BerryKeating1999b}
M.~V.\ Berry, J.~P.\ Keating,
\textit{The Riemann zeros and eigenvalue asymptotics},
SIAM Rev.\ \textbf{41} no.~2 (1999), 236--266.

\bibitem{BerryRobnik1984}
M.~V.\ Berry, M.\ Robnik,
\textit{Semiclassical level spacings when regular and chaotic orbits
coexist},
J.\ Phys.\ A: Math.\ Gen.\ \textbf{17} (1984), 2413--2421.

\bibitem{BrunaultChida2015}
F.\ Brunault, M.\ Chida,
\textit{Regulators for Rankin--Selberg products of modular forms},
Math.\ Res.\ Lett.\ \textbf{23} no.~5 (2016), 1397--1424;
arXiv:1503.04626 [math.NT] (2015).

\bibitem{Castella2024}
F.\ Castella,
\textit{On the Tamagawa number conjecture for CM modular forms},
preprint (2024); arXiv:2407.11891 [math.NT].

\bibitem{Connes1999}
A.\ Connes,
\textit{Trace formula in noncommutative geometry and the zeros of the
Riemann zeta function},
Selecta Math.\ (N.S.) \textbf{5} no.~1 (1999), 29--106;
preprint Bures--IHES M/99/04 (1999).

\bibitem{Cox1989}
D.~A.\ Cox,
\textit{Primes of the Form \(x^2 + ny^2\): Fermat, Class Field Theory,
and Complex Multiplication},
John Wiley \& Sons, New York, 1989.

\bibitem{Damerell1971}
R.~M.\ Damerell,
\textit{L--functions of elliptic curves with complex multiplication
I, II},
Acta Arith.\ \textbf{17} (1970), 287--301; \textbf{19} (1971), 311--317.

\bibitem{Gauss1801}
C.~F.\ Gauss,
\textit{Disquisitiones Arithmeticae},
Leipzig (1801); English translation by A.~A.\ Clarke, Yale University
Press, New Haven, 1966; §§225--237 (genus theory).

\bibitem{Hecke1936}
E.\ Hecke,
\textit{\"Uber Modulfunktionen und die Dirichletschen Reihen mit
Eulerscher Produktentwicklung. I, II},
Math.\ Ann.\ \textbf{114} (1937), 1--28 (orig.\ 1936);
\textbf{114} (1937), 316--351.

\bibitem{IwaniecKowalski2004}
H.\ Iwaniec, E.\ Kowalski,
\textit{Analytic Number Theory},
AMS Colloquium Publications \textbf{53}, American Mathematical Society,
Providence, RI, 2004.

\bibitem{KatzSarnak1999}
N.~M.\ Katz, P.\ Sarnak,
\textit{Zeros of zeta functions and symmetry},
Bull.\ Amer.\ Math.\ Soc.\ (N.S.) \textbf{36} no.~1 (1999), 1--26.
DOI: 10.1090/S0273-0979-99-00766-1.

\bibitem{KeatingSnaith2000}
J.~P.\ Keating, N.~C.\ Snaith,
\textit{Random matrix theory and \(\zeta(1/2 + it)\)},
Commun.\ Math.\ Phys.\ \textbf{214} (2000), 57--89.
DOI: 10.1007/s002200000261.

\bibitem{Mehta2004}
M.~L.\ Mehta,
\textit{Random Matrices},
Pure and Applied Mathematics \textbf{142}, Elsevier/Academic Press,
Amsterdam, 3rd ed., 2004.

\bibitem{Murty2007}
M.~R.\ Murty,
\textit{Problems in Analytic Number Theory},
Graduate Texts in Mathematics \textbf{206}, Springer, New York,
2nd ed., 2007.

\bibitem{Remondiere2026}
K.\ Remondi\`ere,
\textit{The number of \(\Q\)--rational weight--5 CM newforms attached
to an imaginary quadratic field: a theorem (Gauss 1801) and a
conjecture (Rubin 1991)},
companion preprint (2026); to appear, J.\ Number Theory.

\bibitem{Remondiere2026bsd}
K.\ Remondi\`ere,
\textit{Cross--orbit rational ratios of weight--3 CM newforms at
\(h_K = 2\) imaginary quadratic anchors: a \(14\)--anchor empirical
sweep},
companion preprint (2026); in preparation.

\bibitem{Remondiere2026rh}
K.\ Remondi\`ere,
\textit{A per--sector Hilbert--P\'olya conjecture for Hecke
L--functions of imaginary quadratic genus characters: Conjecture
E18--RH--Strong},
preprint (2026); in preparation.

\bibitem{RudnickSarnak1996}
Z.\ Rudnick, P.\ Sarnak,
\textit{Zeros of principal L--functions and random matrix theory},
Duke Math.\ J.\ \textbf{81} no.~2 (1996), 269--322.

\bibitem{Schutt2008}
M.\ Sch\"utt,
\textit{CM newforms with rational coefficients},
Ramanujan J.\ \textbf{19} no.~2 (2009), 187--205;
arXiv:math/0511228 [math.NT] (2008).

\bibitem{Shimura1971}
G.\ Shimura,
\textit{Introduction to the Arithmetic Theory of Automorphic
Functions},
Publications of the Mathematical Society of Japan \textbf{11},
Iwanami Shoten/Princeton University Press, 1971.

\bibitem{Titchmarsh1986}
E.~C.\ Titchmarsh,
\textit{The Theory of the Riemann Zeta--function},
revised by D.~R.\ Heath--Brown, Oxford University Press, Oxford,
2nd ed., 1986.

\end{thebibliography}

\end{document}
